\section{CONCLUSION}

Focusing on the task of selenium nanoparticle (SeNP) cytotoxicity prediction, this study established an integrated research pipeline comprising benchmark dataset model comparison, dedicated selenium nanoparticle modeling, and automated screening of enhanced features. The findings demonstrate that on large-scale heterogeneous tabular nanoparticle datasets, both deep tabular architectures and tree-based models achieve effective predictive performance. However, in selenium nanoparticle scenarios characterized by limited sample sizes and a relatively homogeneous material scope, tree-based models exhibit superior stability and applicability.

Within the dedicated selenium nanoparticle dataset, 383 validated records were curated from 43 peer-reviewed publications spanning the period from 2008 to 2022. While baseline physicochemical features provide preliminary predictive capabilities, the incorporation of material-chemistry enhanced variables—such as elemental Se content, long-term stability, and Se 3d binding energy—substantially improves the fitting capacity of the models. The multi-modal feature fusion framework implemented in this study—concatenation of baseline and enhanced descriptors followed by cross-attention weighting and residual connection—provides a principled approach for integrating heterogeneous feature modalities, enabling the model to dynamically prioritize the most informative descriptors when auxiliary information (such as IC50) is available.

Among the architectures evaluated, XGBoost achieved the highest regression performance under the enhanced feature configuration. However, because this configuration incorporates IC50—which is properly conceptualized as a response-correlated auxiliary variable rather than an a priori physicochemical descriptor—this result is more appropriately designated as an auxiliary benchmark for mechanistic insight and validation, rather than an operational prospective prediction metric. In contrast, Random Forest maintained a balanced compromise among precision, robustness, and interpretability, rendering it a more suitable candidate for subsequent platform deployment and engineering implementation.

Feature selection results further indicate that the optimal combination of enhanced variables does not require the exhaustive set of candidate descriptors; utilizing solely IC50, elemental Se content, long-term stability, and Se 3d binding energy is sufficient to achieve predictive accuracy comparable to that of the complete feature set. Critically, this finding underscores the necessity of distinguishing two classes of enhanced variables: intrinsic material-chemistry descriptors (elemental Se content, stability, and binding energy) that embody transferable and cumulative feature engineering value within a single-material domain, and response-correlated auxiliary variables (IC50) that serve as diagnostic references when historical experimental data are available. This distinction carries direct implications for future data acquisition: experimental design should prioritize measurements of intrinsic chemical attributes that can be leveraged in prospective screening scenarios.

Overall, this study delineates a structured pathway encompassing data curation, model benchmarking, feature screening, and multi-modal fusion to support the development of selenium nanoparticle toxicity prediction platforms. The \textbf{enhanced-feature XGBoost model} provides a strong diagnostic reference ($R^2 = 0.957$), while the material-chemistry-enhanced Random Forest serves as a more deployment-ready baseline. Importantly, the self-synthesized SeNP datasets currently under characterization will enable rigorous external validation of these models in future work, allowing us to assess their generalization capacity in prospective scenarios where IC50 is not available. Collectively, these results establish a foundational framework for computationally guided nanotoxicity assessment and experimental priority screening.
